\section{Clase 15 - Pr\'actica 6 (Resoluci\'on num\'erica de ecuaciones no lineales)}

\subsection{M\'etodo de la secante}

Comenzamos con dos valores $x_{0}$ y $x_{1}$ cercanos a la ra\'iz (podr\'ian estar los dos a la izquierda o a la derecha de la ra\'iz) y construimos $x_{2}$ como la ra\'iz de la recta secante que pasa por $(x_{0}, f(x_{0}))$ y $(x_{1}, f(x_{1}))$.

$$\frac{f(x_{1}) - f(x_{0})}{x_{1} - x_{0}}(x - x_{1}) + f(x_{1}) = 0$$

$$\Rightarrow x_{2} = x_{1} - \frac{f(x_{1})}{\frac{f(x_{1}) - f(x_{0})}{x_{1} - x_{0}}}$$

\begin{itemize}
    \item Tomo $x_{0}$ y $x_{1}$ dos valores aproximados de $r$, la ra\'iz de $f$.\\

    \item Para $n = 1, 2, ..., x_{n + 1} = x_{n} - \frac{f(x_{n})}{\frac{f(x_{n}) - f(x_{n - 1})}{x_{n} - x_{n - 1}}}$
\end{itemize}

\underline{Error en el m\'etodo de la secante}

$$e_{n + 1} = \frac{f''(\eta_{n})}{2f'(\theta_{n})} = e_{n}e_{n - 1}$$

con $\eta_{n}, \theta_{n} \in [a, b]$.\\

\textit{Idea sobre el orden de convergencia}: si $p$ es tal que $e_{n + 1} \sim C|e_{n}|^{p} \Rightarrow |e_{n + 1}| \sim C|e_{n}|^{p} \sim C(C|e_{n - 1}|^{p})^{p} \sim C^{p + 1}|e_{n - 1}|^{p^{2}}$.\\

Por la f\'ormula del error sabemos que $|e_{n + 1}| \sim K|e_{n}||e_{n - 1}|$ y reemplazando en lo anterior nos queda

$$K|e_{n}||e_{n - 1}| \sim C^{p + 1}|e_{n - 1}|^{p^{2}}$$

y de aqu\'i

$$KC|e_{n - 1}|^{p}|e_{n - 1}| \sim C^{p + 1}|e_{n - 1}|^{p^{2}}$$

Entonces

$$\frac{K}{C^{p}}|e_{n - 1}|^{p + 1} \sim |e_{n - 1}|^{p^{2}} \Rightarrow |e_{n - 1}|^{p^{2} - p - 1} \sim cte$$\\

Deber\'ia ser entonces $p^{2} - p - 1 = 0$ y la soluci\'on positiva de esta ecuaci\'on en $p = \frac{1 + \sqrt{5}}{2}$ (el n\'umero de oro).

\subsection{M\'etodo de punto fijo}

La idea ahora es reformular el problema de "hallar $x$ tal que $f(x) = 0$" en otro equivalente de la forma "hallar $x$ tal que $x = g(x)$".\\

\textit{Ejemplo}: $4x^{3} - 2x - 5 = 0$ es equivalente a

\begin{itemize}
    \item $x = \frac{4x^{3} - 5}{2}$, donde $g_{1}(x) = \frac{4x^{3} - 5}{2}$
    
    \item $x = \sqrt[3]{\frac{2x + 5}{4}}$, donde $g_{2}(x) = \sqrt[3]{\frac{2x + 5}{4}}$
    
    \item $x = \frac{2x + 5}{4x^{2}}$, donde $g_{3}(x) = \frac{2x + 5}{4x^{2}}$
\end{itemize}

\textit{Cuidado}: No cualquier reformulaci\'on va a servir.\\

\underline{Definici\'on}: Decimos que $r$ es punto fijo de $g$ si $g(r) = r$.\\

\begin{itemize}
    \item Tomo $x_{0}$ una aproximaci\'on de $r$.

    \item Para $n = 0, 1, 2, ...$, $x_{n + 1} = g(x_{n})$ .
\end{itemize}

\textit{Observaci\'on}: el m\'etodo de Newton-Raphson se puede ver como un m\'etodo de punto fijo para el problema $x = g(x)$ con $g(x) = x - \frac{f(x)}{f'(x)}$\\

¿Bajo que condiciones converge el m\'etodo de punto fijo?\\

\underline{Definici\'on}: Decimos que una funci\'on $g: [a, b] \rightarrow \mathbb{R}$ es contractiva (o es una contracci\'on) en el intervalo $[a, b]$ si existe una constante $L, 0 \leq L < 1$ tal que $|g(x) - g(y)| \leq L|x - y|$ para cualquiera $x, y \in [a, b]$\\

\underline{Teorema}: Sea $g: [a, b] \rightarrow \mathbb{R}$ continua tal que

\begin{itemize}
    \item $g([a, b]) \subseteq [a, b]$
    
    \item $g$ es contractiva en $[a, b]$ con cte $0 \leq L < 1$, entonces el m\'etodo de punto fijo $x_{n + 1} = g(x_{n})$ converge al (\'unico) punto fijo $r \in [a, b]$. Es decir, $\lim_{n \to \infty} x_{n} = r$ con $g(r) = r$.
\end{itemize}

¿C\'omo verifico que $g$ es contractiva?\\

Por el Teorema de Valor Medio de Lagrange, si $g$ es derivable podemos escribir para cualesquiera $x, y \in [a, b]$:\\

$$|g(x) - g(y)| = |g'(\theta)||x - y| \leq \max_{\theta \in [a, b]} |g'(\theta)||x - y|$$

Luego, para ver que $g$ es contractiva, basta con encontrar $0 \leq L < 1$ tal que $\max_{\theta \in [a, b]} |g'(\theta)| \leq L$.\\

¿Por qu\'e converge el m\'etodo?\\

Al pedir $g([a, b]) \subseteq [a, b]$ garantizamos que si $x_{0} \in [a, b]$ entonces toda la sucesi\'on quedar\'a en $[a, b]$:

$$x_{0} \in [a, b] \Rightarrow x_{1} = g(x_{0}) \in [a, b] \Rightarrow x_{2} = g(x_{1}) \in [a, b]\ ...$$

Como $g$ es contractiva,

$$|e_{n + 1}| = |x_{n + 1} - r| = |g(x_{n}) - g(r)| \leq L|x_{n} - r| = L|e_{n}|$$

y con este razonamiento se obtiene

$$|e_{n + 1}| \leq L|e_{n}| \leq L^{2}|e_{n - 1}| \leq ... \leq L^{n + 1}|e_{0}| \rightarrow_{n \to +\infty} 0$$

pues $L < 1$.\\

\underline{Error de punto fijo}

\begin{itemize}
    \item $|x_{n} - r| \leq L^{n}|x_{0} - r|$

    \item $|x_{n} - r| \leq \frac{L^{n}}{1 - L}|x_{1} - x_{0}|$

    \item $|x_{n} - r| \leq \frac{L}{1 - L}|x_{n} - x_{n - 1}|$
\end{itemize}